\section{Proposta de algoritmo para BFS}
Este trabalho propõe a utilização do algoritmo de BFS em fronteiras, discutido na
Seção~\ref{sec:bfs-fronteira}. Também é discutida a paralelização do algoritmo
na Seção~\ref{sec:bfs-paralelizacao}. É proposta uma divisão do grafo entre
workers na Seção~\ref{sec:bfs-divisao-grafo}, é demonstrado um exemplo de BFS
paralela na Seção~\ref{sec:bfs-exemplo} e, por último, é discutida a
implementação do algoritmo em GPUs na Seção~\ref{sec:bfs-gpu}.

\subsection{BFS em Fronteiras}
\label{sec:bfs-fronteira}
O algoritmo de BFS também pode ser implementado como uma descoberta por fronteiras. Inicializa-se a primeira fronteira com o vértice inicial. Em seguida, para cada vértice presente na fronteira atual, examinam-se todos os seus vizinhos; aqueles que ainda não foram visitados são adicionados à próxima fronteira. Esse processo se repete sucessivamente, até que a fronteira atual esteja vazia.

Uma versão simplificada do algoritmo é apresentada a seguir:

\begin{lstlisting}
BFS(G, origem):
    para cada vertice v em G:
        pai[v] = none

    pai[origem] = origem
    fronteira = {origem}

    enquanto fronteira nao esta vazia:
        proxima_fronteira = vazia
        para cada vertice u na fronteira:
            para cada vizinho v de u:
                se pai[v] == none:
                    pai[v] = u
                    adiciona v na proxima_fronteira
        fronteira = proxima_fronteira
\end{lstlisting}

\subsection{Exemplo Visual de Fronteiras}

Considere o grafo da Figura~\ref{fig:bfs-step0}. A busca começa no vértice
\texttt{A}. Em cada passo, os vértices em azul representam a fronteira atual,
os vértices em cinza já foram visitados em níveis anteriores, e os vértices em
branco ainda não foram descobertos. As arestas destacadas em azul formam a
árvore de BFS construída até aquele momento. Essa árvore registra, para cada
vértice descoberto, qual foi o vértice responsável por sua primeira descoberta.
Assim, cada aresta azul liga um vértice ao seu pai na busca, permitindo recuperar
um caminho mínimo, em número de arestas, da origem \texttt{A} até qualquer
vértice alcançável.

\begin{figure}[H]
\centering
\begin{tikzpicture}[scale=1.0]
  \node[frontier] (A) at (0,0) {A};
  \node[undiscovered] (B) at (2,1) {B};
  \node[undiscovered] (C) at (2,-1) {C};
  \node[undiscovered] (D) at (4,1.5) {D};
  \node[undiscovered] (E) at (4,0) {E};
  \node[undiscovered] (F) at (4,-1.5) {F};
  \node[undiscovered] (G) at (6,0) {G};
  \node[undiscovered] (H) at (6,-2) {H};

  \draw[graphedge] (A) -- (B);
  \draw[graphedge] (A) -- (C);
  \draw[graphedge] (B) -- (D);
  \draw[graphedge] (B) -- (E);
  \draw[graphedge] (C) -- (E);
  \draw[graphedge] (C) -- (F);
  \draw[graphedge] (D) -- (G);
  \draw[graphedge] (E) -- (G);
  \draw[graphedge] (F) -- (G);
  \draw[graphedge] (F) -- (H);
\end{tikzpicture}
\caption{Passo 0 da BFS: a fronteira inicial contém apenas a origem \texttt{A}.}
\label{fig:bfs-step0}
\end{figure}

No passo inicial, a distância de \texttt{A} é zero. A primeira expansão examina
as arestas que saem de \texttt{A}. Como \texttt{B} e \texttt{C} ainda não foram
visitados, ambos são descobertos, recebem distância um e passam a formar a
próxima fronteira.

\begin{figure}[H]
\centering
\begin{tikzpicture}[scale=1.0]
  \node[visited] (A) at (0,0) {A};
  \node[frontier] (B) at (2,1) {B};
  \node[frontier] (C) at (2,-1) {C};
  \node[undiscovered] (D) at (4,1.5) {D};
  \node[undiscovered] (E) at (4,0) {E};
  \node[undiscovered] (F) at (4,-1.5) {F};
  \node[undiscovered] (G) at (6,0) {G};
  \node[undiscovered] (H) at (6,-2) {H};

  \draw[treeedge] (A) -- (B);
  \draw[treeedge] (A) -- (C);
  \draw[graphedge] (B) -- (D);
  \draw[graphedge] (B) -- (E);
  \draw[graphedge] (C) -- (E);
  \draw[graphedge] (C) -- (F);
  \draw[graphedge] (D) -- (G);
  \draw[graphedge] (E) -- (G);
  \draw[graphedge] (F) -- (G);
  \draw[graphedge] (F) -- (H);
\end{tikzpicture}
\caption{Passo 1 da BFS: \texttt{B} e \texttt{C} formam a fronteira de nível 1.}
\label{fig:bfs-step1}
\end{figure}

No segundo passo, a BFS expande todos os vértices da fronteira atual. O vértice
\texttt{B} descobre \texttt{D} e \texttt{E}. O vértice \texttt{C} também possui
uma aresta para \texttt{E}, mas \texttt{E} já pode ter sido descoberto por
\texttt{B}; nesse caso, a implementação mantém apenas um pai para \texttt{E}.
O vértice \texttt{C} também descobre \texttt{F}. Assim, a nova fronteira passa a
ser formada por \texttt{D}, \texttt{E} e \texttt{F}, todos com distância dois.

\begin{figure}[H]
\centering
\begin{tikzpicture}[scale=1.0]
  \node[visited] (A) at (0,0) {A};
  \node[visited] (B) at (2,1) {B};
  \node[visited] (C) at (2,-1) {C};
  \node[frontier] (D) at (4,1.5) {D};
  \node[frontier] (E) at (4,0) {E};
  \node[frontier] (F) at (4,-1.5) {F};
  \node[undiscovered] (G) at (6,0) {G};
  \node[undiscovered] (H) at (6,-2) {H};

  \draw[treeedge] (A) -- (B);
  \draw[treeedge] (A) -- (C);
  \draw[treeedge] (B) -- (D);
  \draw[treeedge] (B) -- (E);
  \draw[graphedge] (C) -- (E);
  \draw[treeedge] (C) -- (F);
  \draw[graphedge] (D) -- (G);
  \draw[graphedge] (E) -- (G);
  \draw[graphedge] (F) -- (G);
  \draw[graphedge] (F) -- (H);
\end{tikzpicture}
\caption{Passo 2 da BFS: \texttt{D}, \texttt{E} e \texttt{F} formam a fronteira de nível 2.}
\label{fig:bfs-step2}
\end{figure}

No terceiro passo, a fronteira de nível dois é expandida. Os vértices
\texttt{D}, \texttt{E} e \texttt{F} possuem arestas para \texttt{G}. O primeiro
que conseguir registrar \texttt{G} como descoberto define seu pai na árvore de
BFS. Na figura, \texttt{E} foi escolhido como pai de \texttt{G}, mas
\texttt{D} ou \texttt{F} também poderiam ser pais válidos, pois todos pertencem
ao nível anterior. Além disso, \texttt{F} descobre \texttt{H}, que também passa
a fazer parte da fronteira de nível três.

\begin{figure}[H]
\centering
\begin{tikzpicture}[scale=1.0]
  \node[visited] (A) at (0,0) {A};
  \node[visited] (B) at (2,1) {B};
  \node[visited] (C) at (2,-1) {C};
  \node[visited] (D) at (4,1.5) {D};
  \node[visited] (E) at (4,0) {E};
  \node[visited] (F) at (4,-1.5) {F};
  \node[frontier] (G) at (6,0) {G};
  \node[frontier] (H) at (6,-2) {H};

  \draw[treeedge] (A) -- (B);
  \draw[treeedge] (A) -- (C);
  \draw[treeedge] (B) -- (D);
  \draw[treeedge] (B) -- (E);
  \draw[graphedge] (C) -- (E);
  \draw[treeedge] (C) -- (F);
  \draw[graphedge] (D) -- (G);
  \draw[treeedge] (E) -- (G);
  \draw[graphedge] (F) -- (G);
  \draw[treeedge] (F) -- (H);
\end{tikzpicture}
\caption{Passo 3 da BFS: \texttt{G} e \texttt{H} são descobertos no nível 3.}
\label{fig:bfs-step3}
\end{figure}

Depois da expansão de \texttt{G}, nenhum novo vértice é encontrado. A próxima
fronteira fica vazia e o algoritmo termina. O resultado é uma árvore de BFS na
qual cada vértice alcançável possui um pai em um nível imediatamente anterior ao
seu.

\subsection{Paralelização}
\label{sec:bfs-paralelizacao}
A BFS baseada em fronteiras oferece uma oportunidade natural de paralelismo, todos
os vértices de uma mesma fronteira podem ser processados simultaneamente, e suas
arestas podem ser expandidas em paralelo. O principal cuidado é controlar a
descoberta dos vértices da próxima fronteira, pois múltiplas arestas podem tentar
descobrir o mesmo vértice ao mesmo tempo. Estratégias de paralelismo e controle foram
discutidas em \cite{merrillgpu}. 

\subsection{Divisão do Grafo entre Workers}
\label{sec:bfs-divisao-grafo}
Neste texto, o termo \textit{worker} é usado de forma geral para representar uma
unidade de execução capaz de armazenar parte do grafo e realizar parte do trabalho
da BFS. Um worker pode corresponder, por exemplo, a um processo, uma máquina, uma
GPU, uma CPU ou outro recurso de execução paralelo. Os detalhes de como esses
workers são implementados serão discutidos posteriormente. Essa abstração é
comum em sistemas de processamento paralelo de grafos, nos quais a computação é
dividida entre unidades que armazenam partes do grafo e trocam informações sobre
vértices descobertos durante a execução \cite{merrillgpu}.

Em uma implementação paralela ou distribuída, o grafo não precisa ficar inteiro
em uma única unidade de execução. Uma estratégia comum é dividir seus vértices entre vários
\textit{workers}. Cada worker fica responsável por armazenar e processar uma
parte dos vértices e de suas listas de adjacência. Durante a BFS, quando um
worker expande um vértice local, ele examina suas arestas; se o vizinho pertence
a outro worker, a descoberta desse vizinho precisa ser comunicada ao worker que
o possui \cite{merrillgpu}.

Uma forma simples de particionamento é a distribuição cíclica. Com $n_workers$
workers, o dono de um vértice pode ser calculado por:

\begin{lstlisting}
worker(v) = v mod n_workers
\end{lstlisting}

Assim, os vértices são distribuídos de maneira alternada entre os workers. Essa
estratégia não garante balanceamento perfeito para todos os grafos, mas é simples
e evita concentrar intervalos inteiros de vértices em um único worker. A
Tabela~\ref{tab:worker-partition} mostra um exemplo com 16 vértices.

\begin{table}[H]
\centering
\begin{tabular}{c|c}
Worker & Vértices pertencentes ao worker \\
\hline
0 & 0, 4, 8, 12 \\
1 & 1, 5, 9, 13 \\
2 & 2, 6, 10, 14 \\
3 & 3, 7, 11, 15 \\
\end{tabular}
\caption{Particionamento cíclico de 16 vértices entre 4 workers.}
\label{tab:worker-partition}
\end{table}

A Figura~\ref{fig:grafo-global-workers} mostra a visão global de um grafo com 16
vértices. Essa visão representa o grafo completo antes de considerar qual worker
é responsável por cada vértice.

\begin{figure}[H]
\centering
\begin{tikzpicture}[scale=0.75]
  \foreach \i/\x/\y in {0/0/3,1/2/3,2/4/3,3/6/3,
                        4/0/1,5/2/1,6/4/1,7/6/1,
                        8/0/-1,9/2/-1,10/4/-1,11/6/-1,
                        12/0/-3,13/2/-3,14/4/-3,15/6/-3} {
    \node[vertex] (v\i) at (\x,\y) {\i};
  }

  \draw[graphedge] (v0) -- (v1);
  \draw[graphedge] (v1) -- (v2);
  \draw[graphedge] (v2) -- (v3);
  \draw[graphedge] (v0) -- (v4);
  \draw[graphedge] (v1) -- (v5);
  \draw[graphedge] (v2) -- (v6);
  \draw[graphedge] (v3) -- (v7);
  \draw[graphedge] (v4) -- (v5);
  \draw[graphedge] (v5) -- (v6);
  \draw[graphedge] (v6) -- (v7);
  \draw[graphedge] (v4) -- (v8);
  \draw[graphedge] (v5) -- (v9);
  \draw[graphedge] (v6) -- (v10);
  \draw[graphedge] (v7) -- (v11);
  \draw[graphedge] (v8) -- (v9);
  \draw[graphedge] (v9) -- (v10);
  \draw[graphedge] (v10) -- (v11);
  \draw[graphedge] (v8) -- (v12);
  \draw[graphedge] (v9) -- (v13);
  \draw[graphedge] (v10) -- (v14);
  \draw[graphedge] (v11) -- (v15);
  \draw[graphedge] (v12) -- (v13);
  \draw[graphedge] (v13) -- (v14);
  \draw[graphedge] (v14) -- (v15);
  \draw[graphedge] (v1) -- (v6);
  \draw[graphedge] (v4) -- (v9);
  \draw[graphedge] (v6) -- (v11);
  \draw[graphedge] (v10) -- (v15);
\end{tikzpicture}
\caption{Visão global de um grafo com 16 vértices antes do particionamento.}
\label{fig:grafo-global-workers}
\end{figure}

A Figura~\ref{fig:grafo-particionado-workers} mostra o mesmo grafo particionado
entre quatro workers. As cores indicam o dono de cada vértice. As arestas que
ligam vértices de cores diferentes representam comunicações potenciais durante a
BFS, pois a expansão de um vértice em um worker pode descobrir um vértice que
pertence a outro worker.

\begin{figure}[H]
\centering
\begin{tikzpicture}[scale=0.75]
  \foreach \i/\x/\y/\c in {0/0/3/blue!25,1/2/3/red!25,2/4/3/green!25,3/6/3/orange!35,
                           4/0/1/blue!25,5/2/1/red!25,6/4/1/green!25,7/6/1/orange!35,
                           8/0/-1/blue!25,9/2/-1/red!25,10/4/-1/green!25,11/6/-1/orange!35,
                           12/0/-3/blue!25,13/2/-3/red!25,14/4/-3/green!25,15/6/-3/orange!35} {
    \node[circle, draw, fill=\c, minimum size=8mm, inner sep=0pt] (p\i) at (\x,\y) {\i};
  }

  \draw[graphedge] (p0) -- (p1);
  \draw[graphedge] (p1) -- (p2);
  \draw[graphedge] (p2) -- (p3);
  \draw[graphedge] (p0) -- (p4);
  \draw[graphedge] (p1) -- (p5);
  \draw[graphedge] (p2) -- (p6);
  \draw[graphedge] (p3) -- (p7);
  \draw[graphedge] (p4) -- (p5);
  \draw[graphedge] (p5) -- (p6);
  \draw[graphedge] (p6) -- (p7);
  \draw[graphedge] (p4) -- (p8);
  \draw[graphedge] (p5) -- (p9);
  \draw[graphedge] (p6) -- (p10);
  \draw[graphedge] (p7) -- (p11);
  \draw[graphedge] (p8) -- (p9);
  \draw[graphedge] (p9) -- (p10);
  \draw[graphedge] (p10) -- (p11);
  \draw[graphedge] (p8) -- (p12);
  \draw[graphedge] (p9) -- (p13);
  \draw[graphedge] (p10) -- (p14);
  \draw[graphedge] (p11) -- (p15);
  \draw[graphedge] (p12) -- (p13);
  \draw[graphedge] (p13) -- (p14);
  \draw[graphedge] (p14) -- (p15);
  \draw[graphedge] (p1) -- (p6);
  \draw[graphedge] (p4) -- (p9);
  \draw[graphedge] (p6) -- (p11);
  \draw[graphedge] (p10) -- (p15);

  \node[draw, fill=blue!25, minimum width=7mm, minimum height=4mm] at (8.2,2.5) {};
  \node[anchor=west] at (8.7,2.5) {Worker 0};
  \node[draw, fill=red!25, minimum width=7mm, minimum height=4mm] at (8.2,1.7) {};
  \node[anchor=west] at (8.7,1.7) {Worker 1};
  \node[draw, fill=green!25, minimum width=7mm, minimum height=4mm] at (8.2,0.9) {};
  \node[anchor=west] at (8.7,0.9) {Worker 2};
  \node[draw, fill=orange!35, minimum width=7mm, minimum height=4mm] at (8.2,0.1) {};
  \node[anchor=west] at (8.7,0.1) {Worker 3};
\end{tikzpicture}
\caption{Mesmo grafo particionado entre 4 workers usando a regra $worker(v) = v \bmod 4$.}
\label{fig:grafo-particionado-workers}
\end{figure}

\subsection{Exemplo de BFS Paralela com 4 Workers}
\label{sec:bfs-exemplo}
Considere uma BFS iniciada no vértice 0 do grafo particionado da
Figura~\ref{fig:grafo-particionado-workers}. No passo inicial, apenas o worker 0
possui trabalho, pois ele é o dono do vértice 0. Ao expandir 0, são descobertos
os vértices 1 e 4. O vértice 4 pertence ao próprio worker 0, enquanto o vértice 1
pertence ao worker 1. Portanto, a nova fronteira já fica distribuída entre mais
de um worker.

Após a expansão dos vértices 1 e 4, a BFS chega ao passo mostrado na
Figura~\ref{fig:bfs-paralela-step2}. Nesse momento, a fronteira global de
nível 2 é formada por vértices armazenados em workers diferentes. O worker 0
possui o vértice 8 na fronteira, o worker 1 possui os vértices 5 e 9, o worker 2
possui os vértices 2 e 6, e o worker 3 ainda não possui vértices na fronteira
atual.

A expansão passo a passo é mostrada nas figuras seguintes. Em cada uma delas, a
fronteira global é dividida entre os quatro workers. Os vértices em azul formam
a fronteira local do passo atual, os vértices em cinza já foram visitados, e os
vértices em branco ainda não foram descobertos. Depois que todos os workers
expandem suas fronteiras locais, os vértices recém-descobertos formam a
fronteira do próximo nível.

No passo 0, apenas o worker 0 possui trabalho, pois a origem 0 pertence a ele.

\begin{figure}[H]
\centering
\begin{tikzpicture}[scale=0.72]
  \draw[rounded corners] (-0.8,-0.8) rectangle (5.0,3.4);
  \node[anchor=west] at (-0.6,3.05) {Worker 0};
  \node[frontier] (s0w0v0) at (0,2) {0};
  \node[undiscovered] (s0w0v4) at (2,2) {4};
  \node[undiscovered] (s0w0v8) at (0,0.5) {8};
  \node[undiscovered] (s0w0v12) at (2,0.5) {12};
  \node[anchor=west] at (-0.6,-0.45) {Fronteira local: \{0\}};

  \draw[rounded corners] (6.0,-0.8) rectangle (11.8,3.4);
  \node[anchor=west] at (6.2,3.05) {Worker 1};
  \node[undiscovered] at (6.8,2) {1};
  \node[undiscovered] at (8.8,2) {5};
  \node[undiscovered] at (6.8,0.5) {9};
  \node[undiscovered] at (8.8,0.5) {13};
  \node[anchor=west] at (6.2,-0.45) {Fronteira local: vazia};

  \draw[rounded corners] (-0.8,-6.2) rectangle (5.0,-2.0);
  \node[anchor=west] at (-0.6,-2.35) {Worker 2};
  \node[undiscovered] at (0,-3.4) {2};
  \node[undiscovered] at (2,-3.4) {6};
  \node[undiscovered] at (0,-4.9) {10};
  \node[undiscovered] at (2,-4.9) {14};
  \node[anchor=west] at (-0.6,-5.85) {Fronteira local: vazia};

  \draw[rounded corners] (6.0,-6.2) rectangle (11.8,-2.0);
  \node[anchor=west] at (6.2,-2.35) {Worker 3};
  \node[undiscovered] at (6.8,-3.4) {3};
  \node[undiscovered] at (8.8,-3.4) {7};
  \node[undiscovered] at (6.8,-4.9) {11};
  \node[undiscovered] at (8.8,-4.9) {15};
  \node[anchor=west] at (6.2,-5.85) {Fronteira local: vazia};
\end{tikzpicture}
\caption{Passo 0 da BFS paralela: apenas a origem 0 está na fronteira.}
\label{fig:bfs-paralela-step0}
\end{figure}

Ao expandir 0, são descobertos 1 e 4. O vértice 4 permanece no worker 0, e o
vértice 1 passa a compor a fronteira local do worker 1.

\begin{figure}[H]
\centering
\begin{tikzpicture}[scale=0.72]
  \draw[rounded corners] (-0.8,-0.8) rectangle (5.0,3.4);
  \node[anchor=west] at (-0.6,3.05) {Worker 0};
  \node[visited] at (0,2) {0};
  \node[frontier] at (2,2) {4};
  \node[undiscovered] at (0,0.5) {8};
  \node[undiscovered] at (2,0.5) {12};
  \node[anchor=west] at (-0.6,-0.45) {Fronteira local: \{4\}};

  \draw[rounded corners] (6.0,-0.8) rectangle (11.8,3.4);
  \node[anchor=west] at (6.2,3.05) {Worker 1};
  \node[frontier] at (6.8,2) {1};
  \node[undiscovered] at (8.8,2) {5};
  \node[undiscovered] at (6.8,0.5) {9};
  \node[undiscovered] at (8.8,0.5) {13};
  \node[anchor=west] at (6.2,-0.45) {Fronteira local: \{1\}};

  \draw[rounded corners] (-0.8,-6.2) rectangle (5.0,-2.0);
  \node[anchor=west] at (-0.6,-2.35) {Worker 2};
  \node[undiscovered] at (0,-3.4) {2};
  \node[undiscovered] at (2,-3.4) {6};
  \node[undiscovered] at (0,-4.9) {10};
  \node[undiscovered] at (2,-4.9) {14};
  \node[anchor=west] at (-0.6,-5.85) {Fronteira local: vazia};

  \draw[rounded corners] (6.0,-6.2) rectangle (11.8,-2.0);
  \node[anchor=west] at (6.2,-2.35) {Worker 3};
  \node[undiscovered] at (6.8,-3.4) {3};
  \node[undiscovered] at (8.8,-3.4) {7};
  \node[undiscovered] at (6.8,-4.9) {11};
  \node[undiscovered] at (8.8,-4.9) {15};
  \node[anchor=west] at (6.2,-5.85) {Fronteira local: vazia};
\end{tikzpicture}
\caption{Passo 1 da BFS paralela: os workers 0 e 1 possuem vértices na fronteira.}
\label{fig:bfs-paralela-step1}
\end{figure}

No passo 1, os workers 0 e 1 trabalham ao mesmo tempo: 4 descobre 8 e 9,
enquanto 1 descobre 2, 5 e 6. O vértice 5 também é vizinho de 4, mas basta que
uma das tentativas de descoberta seja aceita para que ele entre na próxima
fronteira. No passo 2, já há trabalho em três workers diferentes.

\begin{figure}[H]
\centering
\begin{tikzpicture}[scale=0.72]
  \draw[rounded corners] (-0.8,-0.8) rectangle (5.0,3.4);
  \node[anchor=west] at (-0.6,3.05) {Worker 0};
  \node[visited] (w0v0) at (0,2) {0};
  \node[visited] (w0v4) at (2,2) {4};
  \node[frontier] (w0v8) at (0,0.5) {8};
  \node[undiscovered] (w0v12) at (2,0.5) {12};
  \node[draw, dashed, circle, minimum size=8mm, inner sep=0pt] (w0g1) at (4,2.8) {1};
  \node[draw, dashed, circle, minimum size=8mm, inner sep=0pt] (w0g9) at (4,0.5) {9};
  \draw[treeedge] (w0v0) -- (w0v4);
  \draw[treeedge] (w0v4) -- (w0v8);
  \draw[graphedge] (w0v8) -- (w0v12);
  \draw[graphedge, dashed] (w0v0) -- (w0g1);
  \draw[graphedge, dashed] (w0v4) -- (w0g9);
  \node[anchor=west] at (-0.6,-0.45) {Fronteira local: \{8\}};

  \draw[rounded corners] (6.0,-0.8) rectangle (11.8,3.4);
  \node[anchor=west] at (6.2,3.05) {Worker 1};
  \node[visited] (w1v1) at (6.8,2) {1};
  \node[frontier] (w1v5) at (8.8,2) {5};
  \node[frontier] (w1v9) at (6.8,0.5) {9};
  \node[undiscovered] (w1v13) at (8.8,0.5) {13};
  \node[draw, dashed, circle, minimum size=8mm, inner sep=0pt] (w1g0) at (10.8,2.8) {0};
  \node[draw, dashed, circle, minimum size=8mm, inner sep=0pt] (w1g6) at (10.8,1.3) {6};
  \draw[treeedge] (w1v1) -- (w1v5);
  \draw[treeedge] (w1v5) -- (w1v9);
  \draw[graphedge] (w1v9) -- (w1v13);
  \draw[graphedge, dashed] (w1v1) -- (w1g0);
  \draw[graphedge, dashed] (w1v1) -- (w1g6);
  \node[anchor=west] at (6.2,-0.45) {Fronteira local: \{5, 9\}};

  \draw[rounded corners] (-0.8,-6.2) rectangle (5.0,-2.0);
  \node[anchor=west] at (-0.6,-2.35) {Worker 2};
  \node[frontier] (w2v2) at (0,-3.4) {2};
  \node[frontier] (w2v6) at (2,-3.4) {6};
  \node[undiscovered] (w2v10) at (0,-4.9) {10};
  \node[undiscovered] (w2v14) at (2,-4.9) {14};
  \node[draw, dashed, circle, minimum size=8mm, inner sep=0pt] (w2g1) at (4,-2.6) {1};
  \node[draw, dashed, circle, minimum size=8mm, inner sep=0pt] (w2g7) at (4,-4.2) {7};
  \draw[treeedge] (w2g1) -- (w2v2);
  \draw[treeedge] (w2g1) -- (w2v6);
  \draw[graphedge] (w2v6) -- (w2v10);
  \draw[graphedge] (w2v10) -- (w2v14);
  \draw[graphedge, dashed] (w2v6) -- (w2g7);
  \node[anchor=west] at (-0.6,-5.85) {Fronteira local: \{2, 6\}};

  \draw[rounded corners] (6.0,-6.2) rectangle (11.8,-2.0);
  \node[anchor=west] at (6.2,-2.35) {Worker 3};
  \node[undiscovered] (w3v3) at (6.8,-3.4) {3};
  \node[undiscovered] (w3v7) at (8.8,-3.4) {7};
  \node[undiscovered] (w3v11) at (6.8,-4.9) {11};
  \node[undiscovered] (w3v15) at (8.8,-4.9) {15};
  \node[draw, dashed, circle, minimum size=8mm, inner sep=0pt] (w3g2) at (10.8,-2.6) {2};
  \node[draw, dashed, circle, minimum size=8mm, inner sep=0pt] (w3g6) at (10.8,-4.2) {6};
  \draw[graphedge] (w3v3) -- (w3v7);
  \draw[graphedge] (w3v7) -- (w3v11);
  \draw[graphedge] (w3v11) -- (w3v15);
  \draw[graphedge, dashed] (w3g2) -- (w3v3);
  \draw[graphedge, dashed] (w3g6) -- (w3v7);
  \draw[graphedge, dashed] (w3g6) -- (w3v11);
  \node[anchor=west] at (6.2,-5.85) {Fronteira local: vazia};
\end{tikzpicture}
\caption{Passo 2 da BFS paralela: a fronteira global está distribuída entre os workers 0, 1 e 2. Vértices tracejados representam vizinhos remotos mantidos apenas para ilustrar comunicações entre workers.}
\label{fig:bfs-paralela-step2}
\end{figure}

No passo 2, a expansão dos vértices 2 e 6 no worker 2 gera trabalho para o
worker 3, pois existem arestas para os vértices 3, 7 e 11. Ao mesmo tempo, os
workers 0, 1 e 2 descobrem outros vértices locais, formando a próxima fronteira.

\begin{figure}[H]
\centering
\begin{tikzpicture}[scale=0.72]
  \draw[rounded corners] (-0.8,-0.8) rectangle (5.0,3.4);
  \node[anchor=west] at (-0.6,3.05) {Worker 0};
  \node[visited] at (0,2) {0};
  \node[visited] at (2,2) {4};
  \node[visited] at (0,0.5) {8};
  \node[frontier] at (2,0.5) {12};
  \node[anchor=west] at (-0.6,-0.45) {Fronteira local: \{12\}};

  \draw[rounded corners] (6.0,-0.8) rectangle (11.8,3.4);
  \node[anchor=west] at (6.2,3.05) {Worker 1};
  \node[visited] at (6.8,2) {1};
  \node[visited] at (8.8,2) {5};
  \node[visited] at (6.8,0.5) {9};
  \node[frontier] at (8.8,0.5) {13};
  \node[anchor=west] at (6.2,-0.45) {Fronteira local: \{13\}};

  \draw[rounded corners] (-0.8,-6.2) rectangle (5.0,-2.0);
  \node[anchor=west] at (-0.6,-2.35) {Worker 2};
  \node[visited] at (0,-3.4) {2};
  \node[visited] at (2,-3.4) {6};
  \node[frontier] at (0,-4.9) {10};
  \node[undiscovered] at (2,-4.9) {14};
  \node[anchor=west] at (-0.6,-5.85) {Fronteira local: \{10\}};

  \draw[rounded corners] (6.0,-6.2) rectangle (11.8,-2.0);
  \node[anchor=west] at (6.2,-2.35) {Worker 3};
  \node[frontier] at (6.8,-3.4) {3};
  \node[frontier] at (8.8,-3.4) {7};
  \node[frontier] at (6.8,-4.9) {11};
  \node[undiscovered] at (8.8,-4.9) {15};
  \node[anchor=west] at (6.2,-5.85) {Fronteira local: \{3, 7, 11\}};
\end{tikzpicture}
\caption{Passo 3 da BFS paralela: todos os workers possuem vértices na fronteira.}
\label{fig:bfs-paralela-step3}
\end{figure}

Após o passo 3, apenas os workers 2 e 3 continuam com novos vértices a expandir.

\begin{figure}[H]
\centering
\begin{tikzpicture}[scale=0.72]
  \draw[rounded corners] (-0.8,-0.8) rectangle (5.0,3.4);
  \node[anchor=west] at (-0.6,3.05) {Worker 0};
  \node[visited] at (0,2) {0};
  \node[visited] at (2,2) {4};
  \node[visited] at (0,0.5) {8};
  \node[visited] at (2,0.5) {12};
  \node[anchor=west] at (-0.6,-0.45) {Fronteira local: vazia};

  \draw[rounded corners] (6.0,-0.8) rectangle (11.8,3.4);
  \node[anchor=west] at (6.2,3.05) {Worker 1};
  \node[visited] at (6.8,2) {1};
  \node[visited] at (8.8,2) {5};
  \node[visited] at (6.8,0.5) {9};
  \node[visited] at (8.8,0.5) {13};
  \node[anchor=west] at (6.2,-0.45) {Fronteira local: vazia};

  \draw[rounded corners] (-0.8,-6.2) rectangle (5.0,-2.0);
  \node[anchor=west] at (-0.6,-2.35) {Worker 2};
  \node[visited] at (0,-3.4) {2};
  \node[visited] at (2,-3.4) {6};
  \node[visited] at (0,-4.9) {10};
  \node[frontier] at (2,-4.9) {14};
  \node[anchor=west] at (-0.6,-5.85) {Fronteira local: \{14\}};

  \draw[rounded corners] (6.0,-6.2) rectangle (11.8,-2.0);
  \node[anchor=west] at (6.2,-2.35) {Worker 3};
  \node[visited] at (6.8,-3.4) {3};
  \node[visited] at (8.8,-3.4) {7};
  \node[visited] at (6.8,-4.9) {11};
  \node[frontier] at (8.8,-4.9) {15};
  \node[anchor=west] at (6.2,-5.85) {Fronteira local: \{15\}};
\end{tikzpicture}
\caption{Passo 4 da BFS paralela: apenas os workers 2 e 3 continuam com fronteiras locais.}
\label{fig:bfs-paralela-step4}
\end{figure}

Por fim, a expansão de 14 e 15 não produz novas descobertas. Nesse ponto, a
próxima fronteira fica vazia e a BFS termina.

\begin{figure}[H]
\centering
\begin{tikzpicture}[scale=0.72]
  \draw[rounded corners] (-0.8,-0.8) rectangle (5.0,3.4);
  \node[anchor=west] at (-0.6,3.05) {Worker 0};
  \node[visited] at (0,2) {0};
  \node[visited] at (2,2) {4};
  \node[visited] at (0,0.5) {8};
  \node[visited] at (2,0.5) {12};
  \node[anchor=west] at (-0.6,-0.45) {Fronteira local: vazia};

  \draw[rounded corners] (6.0,-0.8) rectangle (11.8,3.4);
  \node[anchor=west] at (6.2,3.05) {Worker 1};
  \node[visited] at (6.8,2) {1};
  \node[visited] at (8.8,2) {5};
  \node[visited] at (6.8,0.5) {9};
  \node[visited] at (8.8,0.5) {13};
  \node[anchor=west] at (6.2,-0.45) {Fronteira local: vazia};

  \draw[rounded corners] (-0.8,-6.2) rectangle (5.0,-2.0);
  \node[anchor=west] at (-0.6,-2.35) {Worker 2};
  \node[visited] at (0,-3.4) {2};
  \node[visited] at (2,-3.4) {6};
  \node[visited] at (0,-4.9) {10};
  \node[visited] at (2,-4.9) {14};
  \node[anchor=west] at (-0.6,-5.85) {Fronteira local: vazia};

  \draw[rounded corners] (6.0,-6.2) rectangle (11.8,-2.0);
  \node[anchor=west] at (6.2,-2.35) {Worker 3};
  \node[visited] at (6.8,-3.4) {3};
  \node[visited] at (8.8,-3.4) {7};
  \node[visited] at (6.8,-4.9) {11};
  \node[visited] at (8.8,-4.9) {15};
  \node[anchor=west] at (6.2,-5.85) {Fronteira local: vazia};
\end{tikzpicture}
\caption{Passo 5 da BFS paralela: a fronteira global fica vazia e a busca termina.}
\label{fig:bfs-paralela-step5}
\end{figure}

Essa sequência visual mostra que a fronteira global da BFS é a união das
fronteiras locais de cada worker. Em uma implementação paralela, cada worker
expande apenas os vértices da sua fronteira local. Quando uma aresta aponta para
um vértice remoto, a atualização precisa ser enviada ao worker dono daquele
vértice, que poderá processá-lo no próximo passo. Essa organização preserva a
ideia de avanço por níveis da BFS, mas distribui o trabalho e a comunicação entre
as unidades de execução disponíveis.

\subsection{BFS em GPU}
\label{sec:bfs-gpu}
GPUs modernas são otimizadas para executar operações computacionalmente intensivas 
com muitos grupos de threads em paralelo e dependem de concorrência para esconder 
atrasos de memória. No modelo da NVIDIA, threads são agrupadas em \textit{warps}, 
que executam uma mesma instrução sobre dados diferentes, e em blocos de threads, 
que podem compartilhar memória local no multiprocessador. Esse modelo favorece 
cargas de trabalho com grande quantidade de operações independentes e com acesso 
à memória bem organizado \cite{merrillgpu}.

A BFS é atraente para GPUs porque a expansão de uma fronteira pode expor muito
paralelismo. Em uma mesma iteração, vários vértices da fronteira atual podem ser
expandidos simultaneamente, e as arestas incidentes a esses vértices também podem
ser examinadas em paralelo. Quando a fronteira é grande, há trabalho suficiente
para ocupar muitos multiprocessadores. 

Apesar disso, grafos esparsos tornam a BFS difícil em GPU. O padrão de acesso à
memória depende da estrutura do grafo, não de um arranjo regular como em vetores
ou matrizes densas. Ao expandir um vértice, as threads seguem listas de
adjacência que podem apontar para regiões distantes da memória. Além disso, grafos
reais frequentemente possuem distribuição de grau irregular; alguns vértices têm
poucos vizinhos, enquanto outros concentram muitos. Se uma thread ou um warp recebe 
um vértice de grau muito maior que os demais, ocorre desbalanceamento de carga e 
parte da GPU pode ficar subutilizada \cite{merrillgpu}.

Uma forma simples de paralelizar BFS é fazer cada iteração inspecionar todos os 
vértices ou todas as arestas do grafo, procurando quais deles pertencem ao nível atual. 
Essa estratégia se ajusta facilmente ao modelo de dados paralelos da GPU, mas pode
executar trabalho quadrático no pior caso, pois o mesmo grafo é varrido repetidas
vezes ao longo dos níveis da busca. Para grafos com diâmetro não trivial, esse
custo se torna alto. Uma implementação eficiente deve examinar apenas os vértices
e arestas que pertencem às fronteiras ativas de cada iteração, mantendo custo
linear em relação ao número de vértices e arestas, isto é, $O(|V| + |E|)$
\cite{merrillgpu}.

Do ponto de vista de desempenho, permitir muitas descobertas duplicadas de um mesmo 
vértice pode gerar trabalho redundante, pois o mesmo vértice pode ser inserido mais 
de uma vez na fronteira e sua lista de adjacência pode ser expandida repetidamente. 
Por isso, implementações em GPU costumam usar mecanismos de filtragem, marcação e 
remoção local de duplicatas para reduzir esse efeito \cite{merrillgpu}.

Portanto, a BFS em GPU deve equilibrar três objetivos. Primeiro, preservar a
estrutura por níveis da BFS, garantindo que a próxima fronteira contenha apenas
vértices descobertos a partir da fronteira atual. Segundo, manter eficiência de
trabalho, evitando varrer o grafo inteiro a cada iteração. Terceiro, adaptar a
execução irregular do grafo ao modelo de hardware da GPU, organizando fronteiras,
trabalho e comunicação de forma a reduzir desbalanceamento, contenção e acessos
ineficientes à memória
\cite{merrillgpu}.

\section{Implementação do Graph500 em GPU}
Um dos objetivos deste trabalho foi implementar uma versão em GPU dos kernels do
benchmark Graph500 utilizando CUDA e NVSHMEM. Como discutido na
Seção~\ref{sec:graph500-intro}, o Graph500 avalia sistemas a partir de cargas de
trabalho centradas em grafos, incluindo a construção de uma estrutura interna para
o grafo e a execução de buscas em largura. A implementação desenvolvida neste
trabalho mantém essa organização geral: o Kernel 1 constrói uma representação do
grafo a partir da lista de arestas, enquanto o Kernel 2 executa a BFS a partir
das raízes selecionadas pelo benchmark.

A implementação foi baseada na versão de referência do Graph500, disponível em
\url{https://github.com/graph500/graph500}, mas substitui parte do processamento
por kernels CUDA e operações de comunicação NVSHMEM. A versão desenvolvida neste
trabalho está disponível em
\url{https://github.com/RuanPetrus/graph500/tree/feat/nvshem-implementation}, na
branch \texttt{feat/nvshem-implementation}. O foco da implementação está nos
Kernels 1 e 2 do benchmark. Uma versão do Kernel 3 com NVSHMEM também foi
implementada experimentalmente, mas ela não é o objeto principal desta
monografia, pois o trabalho concentra-se na busca em largura.

O uso de NVSHMEM é motivado pelo modelo apresentado na Seção~\ref{sec:nvshmem}.
Em uma BFS distribuída, a expansão de uma aresta pode descobrir um vértice cujo
estado pertence a outro elemento de processamento. Em uma implementação baseada
apenas em MPI, esse tipo de atualização normalmente precisa ser coordenado pelo
host. Com NVSHMEM, uma thread executando em um kernel CUDA pode emitir operações
remotas diretamente para a memória de outro PE, reduzindo a necessidade de
transferir cada descoberta remota para a CPU antes da comunicação.

\subsection{Modelo de execução}

A implementação utiliza um modelo em que cada processo participante do benchmark
é associado a um PE do NVSHMEM e a uma GPU. Cada PE armazena uma partição do
grafo e executa kernels CUDA para processar os vértices sob sua responsabilidade.
A divisão dos vértices segue a estratégia cíclica discutida na
Seção~\ref{sec:bfs-divisao-grafo}: o dono de um vértice é determinado pelo resto
da divisão entre o identificador global do vértice e o número total de workers.
Com isso, cada PE consegue calcular localmente qual worker é responsável por uma
descoberta produzida durante a BFS.

Esse modelo combina três níveis de execução. A CPU inicializa MPI, CUDA e
NVSHMEM, seleciona o dispositivo associado ao processo e coordena as etapas de
alto nível do benchmark. A GPU executa a expansão das fronteiras da BFS em
paralelo. O NVSHMEM fornece as operações de comunicação entre PEs, permitindo que
uma GPU atualize dados remotos quando uma aresta aponta para um vértice
armazenado por outro worker.

\subsection{Kernel 1}

O Kernel 1 do Graph500 corresponde à construção da estrutura interna do grafo a
partir da lista de arestas gerada pelo benchmark, conforme discutido na
Seção~\ref{sec:graph500-kernel1}. Na implementação desenvolvida, essa etapa
preserva a semântica da versão de referência: a entrada continua sendo uma lista
de tuplas de arestas, e a saída é uma estrutura adequada para executar as buscas
dos kernels seguintes.

Para a execução em GPU, o grafo é convertido para uma representação baseada em
CSR (\textit{Compressed Sparse Row}). Essa representação utiliza um vetor de
deslocamentos para indicar onde começa a lista de adjacência de cada vértice e um
vetor de colunas para armazenar os vizinhos. A escolha por CSR é adequada para a
BFS porque permite localizar rapidamente os vizinhos de um vértice expandido, o
que é essencial durante o processamento das fronteiras.

Após a construção da representação do grafo, os dados necessários para a busca
são alocados ou copiados para a memória da GPU associada ao PE. Cada worker
mantém a parte do grafo correspondente aos seus vértices locais, de acordo com a
distribuição cíclica. Dessa forma, quando o Kernel 2 expande um vértice local, a
lista de adjacência necessária para essa expansão já se encontra disponível na
GPU.

Embora o Kernel 1 não execute a BFS propriamente dita, ele influencia diretamente
o desempenho do Kernel 2. Uma representação inadequada pode aumentar o número de
acessos indiretos à memória, desperdiçar largura de banda ou dificultar o
balanceamento entre threads. Por isso, a etapa de construção do grafo foi mantida
compatível com a especificação do Graph500, mas organizada de modo a permitir que
a expansão das fronteiras ocorra na GPU.

\subsection{Kernel 2}

O Kernel 2 corresponde à execução da busca em largura, descrita na
Seção~\ref{sec:graph500-kernel2}. A implementação utiliza o algoritmo baseado em
fronteiras discutido na Seção~\ref{sec:bfs-fronteira}. Em cada nível da BFS,
existe uma fronteira atual contendo os vértices descobertos no nível anterior.
Esses vértices são expandidos em paralelo, e os vizinhos ainda não visitados
formam a próxima fronteira.

Na versão implementada, cada PE mantém uma fronteira local. A fronteira global da
BFS é a união das fronteiras locais de todos os PEs, como apresentado na
Seção~\ref{sec:bfs-exemplo}. Um kernel CUDA é lançado para expandir os vértices
da fronteira local. Durante essa expansão, as threads percorrem as listas de
adjacência dos vértices ativos e tentam registrar os vizinhos ainda não
descobertos no vetor de predecessores.

Quando o vizinho pertence ao mesmo PE, a atualização é local. Quando o vizinho
pertence a outro PE, a thread utiliza uma operação NVSHMEM para tentar atualizar
o estado remoto desse vértice. Essa atualização precisa ser feita de forma
atômica, pois múltiplas threads, possivelmente em GPUs diferentes, podem tentar
descobrir o mesmo vértice no mesmo nível da BFS. Apenas uma dessas tentativas
deve definir o predecessor do vértice; as demais devem perceber que o vértice já
foi descoberto e não devem inseri-lo novamente como uma nova descoberta válida.

De forma simplificada, a execução do Kernel 2 segue o fluxo abaixo:

\begin{lstlisting}
para cada raiz selecionada pelo Graph500:
    inicializa o vetor de predecessores
    inicializa a fronteira com a raiz

    enquanto a fronteira global nao esta vazia:
        cada PE expande sua fronteira local na GPU
        para cada aresta (u, v) encontrada:
            dono = worker(v)
            se v ainda nao foi descoberto:
                registra u como predecessor de v
                adiciona v na proxima fronteira do dono
        sincroniza os PEs
        fronteira atual = proxima fronteira
\end{lstlisting}

Essa organização preserva a semântica da BFS por níveis: os vértices inseridos na
próxima fronteira foram descobertos exclusivamente a partir da fronteira atual.
Ao mesmo tempo, a expansão das arestas é paralelizada na GPU, explorando o grande
número de operações independentes presentes em fronteiras grandes.

\subsection{Sincronização e comunicação}

A implementação do Kernel 2 utiliza um modelo híbrido entre CPU e GPU. A expansão
de cada fronteira é feita por kernels CUDA, e as atualizações remotas de vértices
são emitidas por operações NVSHMEM diretamente a partir da GPU. Entretanto, a CPU
ainda participa da coordenação entre os níveis da BFS. Ao final de cada nível, é
necessário determinar se a próxima fronteira global está vazia. Caso todos os PEs
estejam sem novos vértices, a BFS termina; caso contrário, a próxima fronteira
passa a ser a fronteira atual do nível seguinte.

Essa sincronização global é necessária para preservar a propriedade de níveis da
BFS e para atender à especificação do Graph500. Porém, ela também representa uma
fonte de custo. Grafos com muitos níveis exigem mais sincronizações, enquanto
grafos com fronteiras grandes tendem a expor mais paralelismo por nível. Assim, o
desempenho da implementação depende tanto do custo de expansão das arestas quanto
do custo de comunicação e sincronização entre PEs.

\subsection{Compatibilidade e limitações}

A implementação procura manter compatibilidade com o fluxo do Graph500 descrito
na Seção~\ref{sec:graph500-fluxo}. A geração do grafo, a construção da estrutura
interna, a execução das buscas e a validação continuam seguindo a organização do
benchmark. O resultado do Kernel 2 também mantém o formato esperado pelo
Graph500: uma árvore de BFS representada por predecessores, na qual a raiz aponta
para si mesma e vértices não alcançáveis permanecem sem predecessor válido,
conforme discutido na Seção~\ref{sec:graph500-validacao}.

Essa versão do Graph500 em GPU serviu como base experimental para a biblioteca
proposta neste trabalho. A partir dela foram identificadas as principais
necessidades de uma implementação distribuída de travessia de grafos em GPU:
alocação de memória acessível por múltiplos PEs, representação distribuída do
grafo, expansão paralela de fronteiras, comunicação remota iniciada no dispositivo
e validação dos resultados produzidos. A seção seguinte apresenta a biblioteca
desenvolvida a partir dessas necessidades.

\section{Biblioteca de Travessia de Grafos Exclusiva em GPU}
\label{sec:biblioteca-gpu}

A segunda etapa do desenvolvimento consistiu em refatorar a implementação
experimental do Graph500 em direção a uma biblioteca de travessia de grafos para
GPU. O objetivo dessa refatoração foi separar as principais responsabilidades do
programa em componentes conceituais mais claros: inicialização do ambiente
paralelo, gerenciamento de memória, geração do grafo, construção da representação
interna, execução da BFS, validação e apoio experimental.

A biblioteca foi projetada para manter os dados na GPU sempre que possível. O
host ainda é responsável por tarefas de controle, como inicialização do ambiente,
seleção de dispositivos, disparo das etapas principais e coleta de métricas. No
entanto, a geração do grafo, a construção da estrutura distribuída e a execução
da BFS são realizadas por kernels CUDA e por operações NVSHMEM. Dessa forma, a
biblioteca reduz a participação da CPU no caminho crítico da busca e se aproxima
do modelo apresentado na Seção~\ref{sec:nvshmem}, no qual a comunicação pode ser
iniciada diretamente pelo dispositivo.

O Graph500 foi mantido como caso de uso principal. Assim, a biblioteca preserva
um fluxo compatível com o benchmark discutido na Seção~\ref{sec:graph500-fluxo}:
geração da lista de arestas, construção da estrutura interna do grafo, seleção de
raízes, execução da BFS, validação e cálculo de métricas.

\subsection{Inicialização e workers}

A execução distribuída é organizada em workers. Neste trabalho, cada worker
representa uma unidade de execução associada a um processo, a um PE do NVSHMEM e
a uma GPU. Durante a inicialização, cada worker identifica sua posição no
conjunto global de processos, seleciona um dispositivo CUDA e participa da
inicialização coletiva do ambiente de comunicação.

\subsection{Arenas simétricas e síncronas}

A biblioteca utiliza arenas como mecanismo principal de alocação de memória. Uma
arena mantém um ponteiro base, a capacidade reservada e a quantidade já utilizada.
Novas alocações são feitas avançando esse ponteiro interno com alinhamento
adequado. Essa estratégia reduz o número de chamadas individuais de alocação e
simplifica o gerenciamento de buffers temporários usados durante a geração do
grafo, conversão para CSR, BFS e validação.

Como o NVSHMEM exige simetria nas alocações, a biblioteca não aloca simplesmente
o tamanho local necessário em cada PE. Antes de uma alocação distribuída, os
workers determinam coletivamente o maior tamanho solicitado. Todos então avançam
suas arenas pelo mesmo tamanho máximo. Isso preserva a simetria da região alocada
e permite que um PE calcule corretamente o endereço remoto correspondente a um
endereço local.

\subsection{Gerador R-MAT}
\label{sec:bib-gerador}

A biblioteca implementa em CUDA o gerador de grafos Kronecker/R-MAT utilizado
pelo Graph500, descrito na Seção~\ref{sec:graph500-geracao-grafo}. A entrada da
geração é definida por dois parâmetros principais: \texttt{SCALE}, que determina
o número global de vértices, e \texttt{edgefactor}, que determina o número global
de arestas. Cada PE recebe um intervalo da lista global de arestas e gera apenas
as arestas pertencentes a esse intervalo.

Para cada aresta, o kernel CUDA avança o estado pseudoaleatório até a
posição correspondente, escolhe recursivamente os quadrantes da matriz de
adjacência conforme os parâmetros iniciadores do R-MAT e aplica a permutação dos
identificadores de vértices. O resultado é armazenado como uma lista distribuída
de arestas, acompanhada dos pesos usados pelas etapas posteriores do benchmark.

A geração é paralelizada por aresta. Cada thread CUDA é responsável por produzir
uma ou mais arestas, permitindo que a etapa de geração explore o paralelismo
disponível na GPU. 

\subsection{Conversão do grafo para CSR}
\label{sec:bib-csr}

Após a geração da lista de arestas, a biblioteca converte o grafo para uma
representação distribuída baseada em CSR. Cada worker armazena os vértices que
lhe pertencem, suas listas de adjacência e os pesos associados às arestas. O
formato é semelhante ao CSR discutido na Seção~\ref{sec:graph500-kernel1}, mas
aplicado a uma distribuição unidimensional cíclica dos vértices.

A conversão ocorre em duas etapas principais. Primeiro, um kernel percorre a
lista de arestas e calcula o grau de cada vértice. Como uma aresta pode conectar
vértices pertencentes a PEs diferentes, o cálculo de grau usa operações atômicas
NVSHMEM para incrementar o contador no PE dono do vértice. Arestas de laço, nas
quais origem e destino são iguais, são ignoradas nessa etapa.

Em seguida, os graus são transformados em deslocamentos para as listas de
adjacência. Depois disso, a biblioteca aloca as listas de adjacência e
executa uma segunda etapa para preenchê-las.

O preenchimento das listas também usa operações NVSHMEM. Para cada direção da
aresta, o kernel identifica o PE dono do vértice de origem, incrementa
atomicamente o deslocamento dentro da linha correspondente e escreve o vizinho no
espaço remoto correspondente. Esse procedimento permite construir a estrutura CSR
distribuída diretamente a partir da GPU, sem primeiro reunir a lista completa de
arestas no host.

\subsection{Kernel 2 exclusivo em GPU}
\label{sec:bib-kernel2}

A principal diferença entre a biblioteca refatorada e a versão híbrida anterior
está na execução do Kernel 2. Na versão descrita na seção anterior, o host ainda
controlava a passagem entre fronteiras da BFS. Na versão refatorada, um kernel
CUDA coletivo mantém o laço principal da BFS dentro da GPU. Assim, a CPU não
precisa lançar um novo kernel a cada nível da busca.

O estado da BFS contém o vetor de predecessores, duas fronteiras, contadores
locais, um contador global de fronteira e variáveis auxiliares usadas para
sincronização no dispositivo. O uso de duas fronteiras permite alternar entre a
fronteira atual e a próxima fronteira sem realocar memória durante a busca.

Dentro do kernel, a raiz é inserida na fronteira inicial pelo PE que possui o
vértice. A cada iteração, os PEs somam suas quantidades locais de vértices ativos
em um contador global armazenado em memória simétrica. Se esse contador for zero,
a BFS termina. Caso contrário, cada PE expande sua fronteira local. A expansão é
organizada para que grupos de threads cooperem no processamento dos vértices
ativos e de suas listas de adjacência, reduzindo a ociosidade causada por
vértices com graus diferentes.

Para cada vizinho encontrado, o kernel calcula o PE dono pela regra de
distribuição cíclica. A descoberta é registrada por meio de uma operação atômica
de comparação e troca sobre o vetor de predecessores do PE dono. Se o valor
anterior indicava que o vértice ainda não havia sido descoberto, a thread que
realizou a atualização foi a primeira a descobrir o vértice. Nesse caso, o
vértice é inserido na próxima fronteira do PE dono por meio de uma operação
atômica de incremento no contador remoto da fronteira.

O kernel também implementa sincronização no dispositivo. Existem barreiras locais
entre blocos do mesmo PE e barreiras globais NVSHMEM entre PEs. Essas barreiras
garantem que todas as atualizações da fronteira atual estejam visíveis antes da
troca entre a fronteira atual e a próxima fronteira. Dessa forma, a biblioteca
preserva a propriedade por níveis da BFS sem retornar ao host entre níveis.

\subsection{Cálculo de distâncias e validação}
\label{sec:bib-validacao}

O resultado principal da BFS é o vetor de predecessores, pois esse é o formato
esperado pelo Kernel 2 do Graph500. Para fins de validação e visualização, a
biblioteca também possui uma rotina que calcula as distâncias a partir da árvore
de predecessores. 

A validação implementada na biblioteca segue a ideia da validação do
Graph500 discutida na Seção~\ref{sec:graph500-validacao}. Em vez de comparar a
BFS com uma árvore fixa, a validação verifica propriedades estruturais do
resultado. Primeiro, confere se os predecessores são válidos, se a raiz aponta
para si mesma e se sua distância é zero. Depois, percorre as arestas do CSR
distribuído para verificar se não existe aresta ligando um vértice visitado a um
vértice não visitado, se as distâncias de vértices adjacentes diferem no máximo
em um nível e se cada predecessor declarado corresponde a uma aresta existente.
